\section{Reading a plot as a scientific statement}
\subsection{Coordinates come before color}
A field is more than a rectangular array. It has a location on a grid, a physical
unit, a domain mask and a saved time. A C-grid velocity component lives on its
own faces; a scalar often lives at cell centers. Earth-frame winds, model-frame
winds and terrain-normal transport components cannot be exchanged solely because
all three have velocity units. Studio's coordinate and vector helpers preserve
these distinctions or reject a view that lacks the required geometry.

Cartesian axes can be translated or nonuniform. Terrain-following levels can
have different physical heights in each column. Cylindrical geometry includes
an angular coordinate. If physical coordinates are unavailable, the application
must identify index units rather than attach metres to an index. The source's
physical-domain bounds determine what may be drawn and reduced; allocated halos
do not become extra atmosphere.

The first two slice indices identify the horizontal directions and the third
identifies the vertical direction in the normalized field representation. A
plan view holds the vertical index fixed. A vertical section holds a horizontal
index fixed. The displayed axes and coordinate metadata, rather than an old
internal variable name, establish whether an axis means Cartesian distance,
angle, height or model level. The plot labelled \textbf{X profile}, for example,
is a Cartesian profile along $x$ at fixed $y$ and $z$; the historical component
filename does not make it a radial average.

\subsection{Slice, vector and composite views}
\textbf{Slice} draws a scalar plane. Changing the slice index requests the
corresponding plane; zoom changes the viewport. For window-capable sources these
operations can request a bounded portion of the source rather than load the
entire three-dimensional field. The source shape, window origin, extent, stride
and aggregation explain what was returned. A lower-resolution display must not
be reported as native merely because its pixels fill the window.

For terrain data, a horizontal model-level plane remains a model-level view:
its height can vary horizontally. The vertical-slice path remaps its one-cell
thick bounded window onto separable physical coordinates before Plotly displays
it. This makes height meaningful on the vertical axis, but it is a documented
remapping operation, not an assertion that the original terrain grid was flat.
Retain the corresponding disclosure when exporting or comparing the view.

\textbf{Vector field} draws arrows for a physically matched component family.
The in-plane components depend on the selected section. Arrow density and scale
affect readability; reducing arrow count does not filter the underlying velocity
field into a new scientific result. Check the frame convention and units before
interpreting arrow direction as Earth-relative motion. A section cannot borrow
an unrelated component solely to make a third arrow coordinate available.

\textbf{Composite} combines a filled scalar plane with contour overlays and
available vector context. Each layer must refer to the same admitted source
plane and frame. Contour thresholds state an explicit condition on values; color
and line width are display choices. If a required overlay is absent or still
loading, a composite cannot truthfully claim that all selected layers have
finished. An older completed view may be retained while its replacement is
prepared, with completion and capture bound to the actual displayed identity.

Deselecting the active field is a valid empty state. It should produce a clear
\emph{No data to plot} state and release unused requests, rather than accuse the
run of failing to export a field the user no longer requested. This distinction
also matters to error handling: an empty selection, a pending read and a genuine
missing export are different conditions.

\subsection{Zoom, native cells and smooth display}
A screen has a finite pixel count; a dataset can have many more cells. Studio
therefore separates the source samples, samples admitted for a particular view,
and the raster used to draw that view. Zooming can refine a bounded preview, but
it cannot recover structures absent from the source. A blocky native-cell plot
at extreme zoom is sometimes an honest view of the model's resolution.

Plotly's heatmap smoothing control changes shading only. It is applied to
regular, finite heatmaps with suitable linear axes. Nonuniform coordinates,
missing values, unsuitable axes and incompatible array shapes retain the
original shading. In particular, smoothing must not bridge a missing-data hole
or reinterpret a stretched coordinate as uniform. The control discloses when
original shading is retained.

For two neighbouring finite samples, linear interpolation has the form
\begin{equation}
 q(\xi)=(1-\xi)q_0+\xi q_1,\qquad 0\leq\xi\leq1.
\end{equation}
This interpolant is bounded by those samples. It can make their transition look
smooth but contributes no new independent measurement between them. Two- and
three-dimensional linear reconstruction apply the corresponding tensor-product
weights. Smoothing a threshold surface is a different operation and belongs to
the disclosed presentation mode described below.

The canvas also follows display pixel density so labels and graphical elements
remain legible on high-density screens. That adjustment is independent of the
simulation's cell spacing. A higher-resolution PNG can improve text and edge
appearance; it cannot turn a coarse numerical field into a fine-grid result.
For publication, preserve units, color limits and sampling disclosure alongside
the image rather than treating pixel count as a scientific resolution measure.

\subsection{Profiles, soundings and hodographs}
\textbf{Vertical} displays a scalar against height at a selected column.
\textbf{X profile} displays a scalar along the Cartesian $x$ direction at fixed
$y$ and $z$. Missing support creates gaps rather than a fabricated continuous
line. Check whether the vertical coordinate is physical height or a declared
model-level coordinate; a level number is not an altitude.

\textbf{Hodograph} plots $u(z)$ against $v(z)$. Its two horizontal axes must use
compatible velocity units and equal geometric scale for direction and turning
to remain interpretable. Height is the parameter along the curve, not a third
horizontal wind component. Storm-relative analysis requires an explicit motion
estimate; translating a hodograph changes its reference frame, not the saved
Earth-relative wind.

\textbf{Skew-T} uses a skewed temperature coordinate and logarithmic pressure
representation. Pressure and either temperature or potential temperature are
required; additional thermodynamic curves require their corresponding inputs.
Do not interpret a missing moisture curve as zero humidity. Sounding plots use
a portrait-friendly layout because forcing the diagram into a very wide, short
pane distorts the geometry that readers use to judge its slopes.

The plotted environmental profile and a parcel-theory interpretation are also
different objects. A sounding's appearance alone does not establish CAPE, parcel
origin, entrainment assumptions or a storm's future development. Retain those
assumptions in the analysis or script used to derive a diagnostic.

\subsection{Time series, time-height plots and swaths}
\textbf{Time-series} applies a declared reducer to a field over completed saved
frames. Mean, minimum, maximum and other offered reductions answer different
questions. A native-cell arithmetic mean is not automatically a mass-weighted
or volume-weighted integral, especially on stretched grids. Read the selected
reducer and geometry rather than infer a conservation statement from a smooth
curve. The horizontal coordinate is saved simulation time, not the wall-clock
time required to load the frame.

\textbf{Time-height} retains a vertical level while reducing the other
directions, producing a Hovmöller-style diagram. Irregular output times and
physical height coordinates must remain visible in its geometry. Large gaps
between saved frames are not observations of what happened between them.
Height-bin averaging and native model-level reduction should not be confused.

\textbf{Swath} accumulates a selected per-cell maximum or minimum across saved
frames. This is a history product: an intense swath location need not correspond
to the currently selected frame. Moving-domain data require explicit Earth-frame
placement if the intended result is a geographic storm track. A moving grid's
array index does not stay over the same Earth location. The swath implementation
tracks source motion where the necessary metadata are available.

These views can read many frames. Their cost may therefore be dominated by
storage and reductions rather than graphics. Changing the requested history
range or field should cancel obsolete work, while shared useful reads may stay
active for other consumers. An analysis assembled from fewer saved frames has
less temporal information even if its display interpolates smoothly.

\subsection{Radar products}
The Radar category uses product presets rather than a separate universal radar
solver. It selects supported reflectivity-related views from the available
dataset. A simulated reflectivity diagnostic is not itself an observed radar
scan, and a color scale does not supply beam geometry or observational effects
that the diagnostic did not calculate.

Logarithmic reflectivity requires care when reducing it. Averaging decibels and
averaging linear reflectivity are different operations. For values $D_i$ in dBZ,
a linear-reflectivity mean is
\begin{equation}
 \overline D=10\log_{10}\!\left(\frac{1}{n}
       \sum_{i=1}^{n}10^{D_i/10}\right).
\end{equation}
The sampling disclosure distinguishes a bounded linear-$Z$ mean from ordinary
cell means and extrema. A radar floor in a three-dimensional view is an
additional layer with its own source, reduction and color mapping; it must
belong to the same displayed frame as the volume above it.

\section{Three-dimensional views}
\subsection{Scientific mode and presentation mode}
Three-dimensional views distinguish \textbf{scientific} and
\textbf{presentation} fidelity. Scientific mode uses resolved source data with
declared sampling, interpolation and feature extraction. It does not mean that
every source cell is always resident or every pixel is an exact source sample.
Bounded subsampling, cell means, extrema reductions and linear reconstruction
have distinct provenance because they change what one can infer about peaks.

Presentation mode can add explicitly disclosed surface smoothing or
support-constrained subgrid detail. These operations can help explain a scene's
shape and continuity. They do not establish new resolved vortices, cloud scales
or physical gradients. Changing fidelity is separate from choosing graphics
quality, ray-step density or screen resolution. A lower-cost scientific render
can still be scientifically honest when its reduction is accurately reported.

Sampling provenance records the source lattice, retained lattice and display
lattice independently. It also records coordinate resampling, including
linear-to-uniform resampling of stretched axes or physical-volume terrain
remapping. In a multi-layer scene, the overall disclosure must account for its
coarsest contributing sampling; a native-resolution foreground cannot make a
coarse background native.

\subsection{Volume rendering}
\textbf{Volume} displays a three-dimensional scalar through a color transfer
function and an opacity transfer function. Color answers what value was sampled;
opacity determines how strongly it contributes along the viewing ray. A bright
or solid-looking region is therefore a consequence of both the field and the
chosen transfer function. Different camera directions can expose or hide the
same interior structure without changing the source.

The renderer samples a ray through a bounded volume texture, respecting validity
support and the scene's coordinate transform. Missing cells should not become
opaque zero-valued material. Interpolating near missing support needs validity
handling, not just a color for NaN. Texture capacity and rendering budgets can
require a bounded representation; source disclosure remains attached to it.

Color range, opacity range and any threshold answer different questions. Fixing
a color range across frames makes visual comparison easier; recomputing a range
for each frame can reveal weak structure but can also make unchanged colors
represent different values. Read the colorbar and range policy when interpreting
growth or decay. Outlier-resistant display ranges can improve visibility while
leaving extrema outside the emphasized range; they do not redefine those extrema.

\subsection{Isosurfaces}
\textbf{Isosurface} extracts the set where a scalar crosses a chosen threshold,
using a marching-cubes family of extraction paths. A threshold surface is a
geometric representation of a level set, not a material boundary unless the
scientific interpretation justifies that identification. Its shape depends on
the threshold, source sampling, interpolation and any declared smoothing.

The preparation pipeline can keep scalar backing and extracted geometry in
workers for reuse. Reusing a previous contour is valid only when the current
prepared dataset identity matches the backing that produced it. Equal field
names, equal dimensions or an apparently familiar frame do not prove that a
worker still holds the same bytes. When a source is evicted and prepared again,
its new resident identity forces the appropriate rebuild.

Lighting changes should not require rebuilding the level set. They update its
material and then the committed appearance descriptor, preserving the source
and geometry identity. The capture metadata must describe the lighting actually
drawn. An export containing the new pixels with the old shadow or fill settings
would be an inconsistent scientific figure recipe.

\subsection{Streamlines and parcel paths}
\textbf{Streamlines} integrate through a vector field at one saved time. The
seed region, direction, integration limits and available support determine the
curves. A streamline is tangent to the instantaneous vector field; it is not
automatically the trajectory of an air parcel through a time-varying storm.
Velocity and vorticity streamline choices require the corresponding coherent
vector field and valid coordinate interpretation.

Geometry preparation, vector sampling and seed selection are distinct from
tube width, color and lighting. A thinner line is not a weaker wind. Cropping
the seed region can reveal the requested feature and reduce work, but omitted
seeds do not establish that other flow paths do not exist. Invalid or unsupported
locations terminate or exclude integration according to the extraction contract.

\textbf{3D paths} reads saved Lagrangian trajectory products. It therefore needs
trajectory output, not merely a set of instantaneous wind fields. The companion
\textbf{Vorticity budget} view uses the saved parcel tendency terms. These are
history-dependent products and are not per-frame volume renderers. Missing
trajectory export should disable the feature clearly rather than silently
substitute a streamline.

\subsection{Rendered supercell and coordinated layers}
\textbf{Rendered supercell} combines condensate volume with available layers
such as a reflectivity floor, updraft structure and other offered diagnostics.
Its availability starts with a supported cloud or total-condensate field;
optional layers have their own input requirements. This is a composite
visualization, not an additional microphysics scheme or a reconstruction of
unexported fields.

Shared volume composition coordinates geometry, textures and shader evaluation
across participating layers. Overlapping volumes must use consistent domain
transforms and physical sampling. Independently drawing translucent actors in
an arbitrary order can change the image; the composition path therefore owns
their joint ordering and framebuffer use. Nested outputs additionally need
parent/child spatial ownership so overlapping refinement regions are not counted
twice as visual material.

Optional cloud presentation detail and motion cues are declared display
operations. They must remain supported by the source cloud and its coordinate
frame. A moving presentation is not a newly integrated cloud field. Use
scientific mode and a direct slice when a conclusion depends on whether a
small structure actually exists in the exported model samples.

\subsection{Camera, orientation and appearance controls}
The camera has a position, target and orientation relative to the physical
scene. The compass, domain box, axes and elevation controls help make this
relationship explicit. Physical elevation controls need the loaded grid; a
placeholder range before geometry is available must not be interpreted as a
valid model height.

Color, background, light direction, shadow depth, fill tint and material
properties control appearance. They can improve shape perception but also
emphasize features for purely optical reasons. Record them when comparing
images. In particular, a new lighting setting must update the appearance
receipt after a successful redraw without falsely changing the scientific
frame or geometry that the scene contains.

Shared and independent control scopes allow several panes to coordinate some
choices while retaining their own views. Camera synchronization does not prove
that two datasets align scientifically. Confirm source time, domain, units and
sampling before using synchronized views as a comparison.
